\subsection{Small-modulus prime-phase regression test}
\label{subsec:numerical-prime-resonance}

For the toy interval $2000<q\le4000$, there are exactly
\begin{equation}
 \pi(4000)-\pi(2000)=247
 \label{eq:toy-prime-count}
\end{equation}
primes.  Every one satisfies $q^2\equiv1\pmod{24}$.  Hence, for $\alpha=m/12$,
\begin{equation}
 S_{\mathrm{prime}}(\alpha)
 :=\sum_{2000<q\le4000\atop q\ {\rm prime}}
 e\!\left(\alpha\frac{q^2}{2}\right)
 =247e(m/24),
 \label{eq:toy-prime-resonance}
\end{equation}
and therefore $|S_{\mathrm{prime}}(\alpha)|=247$ exactly.

\begin{table}[htbp]
\centering
\caption{Exact prime-phase coherence on the tested $2$--$3$ resonance lattice.}
\label{tab:toy-prime-resonance}
\small
\begin{tabular}{ccc}
\toprule
$\alpha$ & common phase & $|S_{\mathrm{prime}}(\alpha)|$\\
\midrule
$0$ & $1$ & $247$\\
$1/2$ & $e(1/4)$ & $247$\\
$1/3$ & $e(1/6)$ & $247$\\
$2/3$ & $e(1/3)$ & $247$\\
$1/4$ & $e(1/8)$ & $247$\\
$3/4$ & $e(3/8)$ & $247$\\
$1/6$ & $e(1/12)$ & $247$\\
$5/6$ & $e(5/12)$ & $247$\\
\bottomrule
\end{tabular}
\end{table}

The old modulus-$r$ expression fails the regression test at $a=1,r=3$:
\begin{equation}
 \sum_{u\bmod3\atop(u,3)=1}e(u^2/6)=e(1/6)+e(4/6)=0.
 \label{eq:old-r3-regression-failure}
\end{equation}
The corrected units-modulo-$6$ factor is
\begin{equation}
 \frac1{\varphi(6)}
 \sum_{u\bmod6\atop(u,6)=1}e(u^2/6)
 =e(1/6),
 \label{eq:correct-r3-regression}
\end{equation}
with unit magnitude.

The elevated value at $\alpha=1/20$ is also arithmetic rather than a numerical artifact.  For primes $q>5$, $q^2\equiv1$ or $9\pmod{40}$.  In this interval the class counts are $123$ and $124$, respectively, so
\begin{equation}
 \left|123e(1/40)+124e(9/40)\right|
 =199.828062\ldots.
 \label{eq:one-twentieth-secondary-resonance}
\end{equation}
This illustrates why the complete reduced-square distribution modulo $2r$, not only the denominator $r$, belongs in the rational major-arc model.
